% Anhang zur Diplomarbeit:
\section{Anhang}
\renewcommand{\arraystretch}{1.15}
\begin{longtable}{|>{\raggedright\arraybackslash}p{2.5cm}|>{\raggedright\arraybackslash}p{10.0cm}|>{\raggedleft\arraybackslash}p{2.2cm}|}
\hline
\textbf{Datum}  & \textbf{Tätigkeit} & \textbf{Zeitaufwand in h} \\
\hline
\endfirsthead
\hline
\textbf{Datum}  & \textbf{Tätigkeit} & \textbf{Zeitaufwand in h} \\
\hline
\endhead
\hline
\endfoot

11.04.2025 & erste Besprechung - aktuellen Stand erklären & 1 \\
30.06.2025 & altes Repository durchsehen & 1 \\
30.06.2025 & zweite Besprechung - was gemacht werden muss & 1 \\
30.06.2025 & Diplomarbeitsantrag anfangen & 1 \\
02.07.2025 & dritte Besprechung - Einzelheiten klären + Fragen beantworten & 0.5 \\
02.07.2025 & Raspi flashen + ssh aufsetzen & 2 \\
12.07.2025 & Diplomarbeitsantrag weiterarbeiten & 2 \\
18.07.2025 & erste Version des Diplomarbeitsantrag finalisieren & 2 \\
29.07.2025 & Pseudo MQTT Client aufsetzen & 3 \\
29.07.2025 & Google Drive aufsetzten + Besprechungsprotokoll & 1 \\
30.07.2025 & Grundkonzept erstellen & 2 \\
30.07.2025 & Informationsaustausch & 1 \\
04.08.2025 & Docker Container erstellen & 2 \\
05.08.2025 & Docker Container erstellen & 2 \\
05.08.2025 & Standort des MQtt-Server hinterfragen & 1 \\
05.08.2025 & Gescheiteter Migrationsversuch von Google Drive -> OneDrive & 1 \\
12.08.2025 & MQtt Worker Container erstellen & 2 \\
10.09.2025 & Besprechung über aktuellen Stand & 0.5 \\
24.09.2025 & Diplomarbeitantrag überarbeiten & 0.5 \\
25.09.2025 & Besprechung über Pflichtenheft & 0.5 \\
27.09.2025 & Diplomarbeitsantrag überarbeiten & 2 \\
07.10.2025 & Informieren über FastAPI-MQtt & 0.5 \\
07.10.2025 & Informieren über Redis in-memory-db & 1 \\
12.10.2025 & MQtt Worker Container überarbeiten nach Besprechung & 2 \\
12.10.2025 & Besprechung mit Leitner Lukas über Coolify & 1 \\
12.10.2025 & Testen des MQtt Worker Containers & 1 \\
14.10.2025 & ``Verschönerung'' der Implementierung des MQtt Worker & 1 \\
15.10.2025 & Meeting of Minutes vorbereiten & 0.5 \\
06.11.2025 & Desktop App Backend & 4 \\
07.11.2025 & Besprechung aktueller Stand & 0.5 \\
12.11.2025 & Gedanken über Frontend der Desktop App machen & 1 \\
13.11.2025 & Desktop App Aufbau überdenken & 2 \\
14.11.2025 & Interne Besprechung über Datenbank Schema & 1 \\
15.11.2025 & Desktop App Backend ändern & 2 \\
19.11.2025 & Backend Routen + CRUD hinzufügen/finalisieren & 3 \\
19.11.2025 & Besprechung über aktuellen Stand + Besprechungsnotizen bearbeiten & 1 \\
06.12.2025 & Frontend anfagen & 7 \\
08.12.2025 & Frontend weiterarbeiten & 5 \\
13.12.2025 & Zwischenpräsentation vorbereiten & 4 \\
14.12.2025 & Zwischenpräsentation Plakat erstellen & 2 \\
15.12.2025 & Zwischenpräsentation Plakat + PowerPoint überarbeiten & 4 \\
26.12.2025 & Frontend weiterarbeiten & 5 \\
17.01.2026 & Frontend fast abschließen. Zeitzonen fehler muss noch fixed werden & 5 \\
19.01.2026 & Frontend refactoring & 5 \\
24.01.2026 & Frontend Fehlerbehebung (Zeitzonen gefixed) + Logbuch & 4 \\
30.01.2026 & Frontend nach Besprechung anpassen und teilweise Zusammenführung mit ZEC-API & 5 \\
02.02.2026 & Frontend weiterarbeiten. Manueller Start + Ende -> Manuelle Zeit & 4 \\
06.02.2026 & Frontend fertigstellen & 6 \\
12.02.2026 & MQtt Worker überarbeiten + Frontend Container ändern & 3 \\
12.02.2026 & Metting aktueller Stand + Zusammenführung + Ablauf attempt creation & 0.5 \\
12.02.2026 & Inhaltsverzeichnis für Dokumentation erstellen & 1 \\
16.02.2026 & Versuch Lichtschranke zum laufen zu bringen & 3 \\
18.02.2026 & MQtt Broker Auth hinzufügen + Automatischen builden und uploaden von Container Images zu ghcr + Kubernetes initiales Setup & 6 \\
19.02.2026 & Über Kubernetes informieren. Tutorials schauen & 3 \\
20.02.2026 & Kubernetes weiterarbeiten. Deployment + Service fertig, Ingress eventuell anpassen & 6 \\
03.03.2026 & Frontend Layout ändern & 4 \\
03.03.2026 & Inhaltsverzeichnis überarbeiten + neue DA Vorlage aufsetzen & 3 \\
04.03.2026 & Automatischs builden und uploaden (in ghcr) von jedem Container außer Frontends. Kubernetes überarbeitet & 3 \\
09.03.2026 & Kubernetes fixen, deployment namen fixen, startups fixen & 3 \\
11.03.2026 & Datenbanken in Kubernetes Cluster verschieben & 5 \\
19.03.2026 & Ingres fixen, startups überarbeiten & 4 \\
20.03.2026 & Versuch, mit Desktop App auf Cluster zugreifen & 4 \\
23.06.2026 & Finalisierung der Desktop-App mit automatischen builden von Image + an Kubernetes Cluster weiterarbeiten & 2 \\
07.07.2026 & Finalisierung des Kubernetes Clusters & 3 \\
31.08.2026 & Schreiben der Diplomarbeit & 2 \\
01.09.2026 & Schreiben der Diplomarbeit & 5.5 \\
09.09.2026 & Schreiben der Diplomarbeit & 5 \\
11.09.2026 & Schreiben der Diplomarbeit & 7 \\
12.09.2026 & Schreiben der Diplomarbeit & 3 \\
13.09.2026 & Schreiben der Diplomarbeit & 3 \\

\hline
\textbf{Gesamt} & & \textbf{185.5} \\
\hline
\end{longtable}
\newpage

\section{Verzeichnisse}
% Tabellenverzeichnis:
\subsection{Abbildungsverzeichnis}
{%
\makeatletter
\@starttoc{lof}
\makeatother
}

\newpage
\subsection{Abkürzungsverzeichnis}
{%
\renewcommand*{\glossarysection}[2][]{}
\printnoidxglossary[type=\acronymtype, sort=letter]
}
\newpage
\subsection{Literaturverzeichnis}
\printbibliography[heading=none]

\newpage
\section{Code Beispiele}
\subsection{MQTT-Worker - Erhalt der Nachricht Code Beispiel}

\begin{CodeBox}
def extract_payload(payload):
    mac_address = payload.get("esp_id")
    timestamp = payload.get("timestamp")

    # Clean values
    mac_address = mac_address.split("-", 1)[1]

    add_timestamp(mac_address, timestamp)

    return mac_address, timestamp

def on_connect(client, userdata, flags, rc):
    print(f"Connect with result code {rc}")

def on_message(client, userdata, msg):
    try:
        mac, timestamp = extract_payload(json.loads(msg.payload.decode("utf-8")))
        print(f"Received mac-address: {mac} and timestamp: {timestamp}")
    except json.JSONDecodeError:
        print("Received message is not valid JSON", file=sys.stdout)
\end{CodeBox}

\newpage

\subsection{MQTT-Worker - API Endpoint Code Beispiel}

\begin{CodeBox}
def get_timestamps(mac: str):
    key = f"timestamps:{mac}"
    now = time.time()

    redis_connection.zremrangebyscore(key, 0, now)

    if not redis_connection.exists(key):
        raise MacNotFound(
            message=(
                "No ESP with this MAC address exists or no timestamps yet exist. "
                "If you think it should exist, please send new MQTT data."
            ),
            error_code=404,
        )
    return redis_connection.zrange(key, 0, -1)

@router.get("/{esp_mac}", response_model=Timestamp, status_code=status.HTTP_200_OK)
def get_timestamp_from_esp(esp_mac: str):
    clean_mac = esp_mac.replace("-", ":")
    try:
        timestamps = get_timestamps(clean_mac)
    except MacNotFound as e:
        raise HTTPException(
                status_code=status.HTTP_404_NOT_FOUND,
                detail=(
                    "No ESP with this MAC address exists or no timestamps yet exist. "
                    "If you think it should exist, please send new MQTT data."
                )
            )
    return Timestamp(timestamp=timestamps)
\end{CodeBox}

\newpage

\subsection{Zeitnehmerapplikation - API-Key Authentifizierung}

\begin{CodeBox}
const fetchTeams = async () => {
    try {
      const response = await axios.get<Team[]>(
        `${SERVER_API_URL}/teams/`,
        {
          headers: {
            "x-api-key": API_KEY
          }
        }
      );

      setTeams(response.data);
      if (response.data.length > 0) setSelectedTeam(response.data[0]);
    } catch (error) {
      console.error("Failed to fetch teams", error);
    }
};
\end{CodeBox}

\newpage

\subsection{Zeitnehmerapplikation - Kommunikation mit MQTT-Worker}

\begin{CodeBox}
const fetchStartTimestamps = async () => {
    if (!selectedChallenge) return;
        try {
        const start_1 = await axios.get<{ timestamp: string[] }>(`${MQTT_WORKER_API_URL}/timestamps/${selectedChallenge.esp_mac_start1.replace(/:/g, '-')}`);
        const start_2 = await axios.get<{ timestamp: string[] }>(`${MQTT_WORKER_API_URL}/timestamps/${selectedChallenge.esp_mac_start2.replace(/:/g, '-')}`);
        const combined = [...(start_1.data.timestamp || []), ...(start_2.data.timestamp || [])];
        setStartTimestamps(combined);
    } catch (error) {
      console.error("Failed to fetch start timestamps", error);
    }
};
\end{CodeBox}
